\chapter*{L}
\label{chap:l}

\term{Lagrange's Theorem}
For a finite group $G$ and subgroup $H$, the group is partitioned into cosets of $H$, so $|G|=[G:H]|H|$. Thus the order of $H$ divides the order of $G$. \par\noindent\textbf{Applications:} the theorem restricts possible subgroup orders and supports the study of symmetries, permutation groups, and modular arithmetic.

\term{Lagrange Multipliers}
For an equality constraint $g(x)=c$, a constrained extremum of $f$ often satisfies $\nabla f=\lambda\nabla g$, together with $g(x)=c$. The gradients are parallel because the objective cannot change to first order along the constraint surface. \par\noindent\textbf{Applications:} the method supports constrained design, economics, mechanics, and statistical estimation.

\term{Lagrange Interpolation}
Given distinct nodes $x_i$ and values $y_i$, Lagrange interpolation constructs the unique polynomial of degree at most $n-1$ through the data using $p(x)=\sum_i y_iL_i(x)$, where $L_i(x_j)=\delta_{ij}$. \par\noindent\textbf{Applications:} it estimates values from tables and measurements and is used in numerical integration and scientific computing.

\term{Lagrange Inversion Theorem}
The Lagrange inversion theorem gives coefficients of the local inverse of an analytic function. If $w=z\phi(w)$ with $\phi(0)\ne0$, then coefficients of $w(z)$ can be computed from powers of $\phi$. \par\noindent\textbf{Applications:} it derives generating functions, counts combinatorial structures, and analyzes implicit equations.

\term{Lagrangian (Mechanics)}
The Lagrangian is often $L=T-V$, kinetic minus potential energy. The motion follows the Euler--Lagrange equations $\frac d{dt}(\partial L/\partial\dot q)-\partial L/\partial q=0$. The formulation handles generalized coordinates and constraints naturally. \par\noindent\textbf{Applications:} it models mechanics, optics, field theory, robotics, and control.

\term{Lambda Calculus}
Lambda calculus is a formal language for computation built from variables, function abstraction, and function application, with substitution as its main operation. It separates the idea of a function from a particular programming syntax. \par\noindent\textbf{Applications:} it underlies functional programming, type theory, interpreters, and the mathematical study of computability.

\term{Lambda-term}
A lambda term is a variable, an abstraction $\lambda x.M$ representing a function, or an application $(MN)$ applying one term to another. Bound variables can be renamed, and substitution models computation. \par\noindent\textbf{Applications:} lambda terms describe programs, symbolic reduction, and formal proofs in logic and type systems.

\term{Langlands Program}
A vast network of conjectures, initiated by Robert Langlands in 1967, connecting automorphic forms and representations of reductive groups over global fields to Galois representations and motives. The central reciprocity conjecture, the Langlands correspondence, predicts that $L$-functions of automorphic representations match those of Galois representations, and functoriality asserts functorial transfer between reductive groups. Its local version concerns representations of $p$-adic groups, and it generalizes Artin's conjecture on holomorphy, providing the framework for modularity results such as the proof of Fermat's Last Theorem.

\term{Laplace Equation}
The equation $\Delta u = 0$ for the Laplacian $\Delta$, whose solutions are the harmonic functions. Harmonic functions satisfy the maximum principle and the mean value property and are real analytic. Boundary value problems for the equation prescribe the values of $u$ or of its normal derivative on the boundary, the Dirichlet and Neumann problems; the inhomogeneous Poisson equation $\Delta u = f$ is solved using fundamental solutions and Green's functions. The Laplace equation is the model elliptic equation, governing electrostatics, steady heat flow, and potential flow.

\term{Laplace Transform}
The Laplace transform is $\mathcal L\{f\}(s)=\int_0^\infty e^{-st}f(t)\,dt$. Differentiation becomes multiplication by $s$ together with initial-value terms, so differential equations become algebraic equations in transform space. \par\noindent\textbf{Applications:} it solves ODEs for circuits, control systems, mechanical vibrations, and transient responses.

\term{Laplacian}
The Laplacian is the divergence of the gradient: $\Delta f=\nabla\cdot\nabla f=\sum_i\partial^2f/\partial x_i^2$. It measures local curvature or imbalance of flux. \par\noindent\textbf{Applications:} it governs heat and diffusion, electrostatic potential, wave models, image smoothing, and graph-based learning.

\term{Large Cardinals}
Very large infinite cardinals whose existence cannot be proved in ZFC and whose consistency implies the consistency of ZFC, arranged in a hierarchy of consistency strength. They include inaccessible cardinals, measurable cardinals, Woodin cardinals, and supercompact cardinals. Their theory measures the strength of set-theoretic statements, interacts with the axiom of determinacy and with forcing, and features in inner model theory, where the existence of such cardinals reflects deep structure of the universe of sets beyond the usual cumulative hierarchy.

\term{Large Deviations Principle}
The asymptotic description of rare events: a sequence of probability measures $P_n$ satisfies a large deviations principle with rate function $I$ if $P_n(A) \approx \exp(-n \inf_{x \in A} I(x))$ for suitable sets $A$, where $I \geq 0$ is lower semicontinuous with compact sublevel sets. Cramér's theorem gives the rate function for empirical means of i.i.d. variables, and Sanov's theorem for empirical measures. Varadhan's lemma computes limits of normalized logarithmic expectations. The theory underpins statistics and statistical mechanics.

\term{Lattice (Algebra)}
A lattice is a partially ordered set in which every pair has a least upper bound $a\vee b$ and greatest lower bound $a\wedge b$. Distributive and modular lattices are important special cases. \par\noindent\textbf{Applications:} lattices organize subspaces, ideals, logical propositions, and data or information hierarchies.

\term{Lattice (Geometry)}
A geometric lattice is a discrete subgroup $\Lambda=\{\sum_i m_i\mathbf b_i:m_i\in\mathbb Z\}$ of $\mathbb R^n$. A basis generates a fundamental parallelepiped whose volume is $|\det(\mathbf b_1,\ldots,\mathbf b_n)|$ in full rank. \par\noindent\textbf{Applications:} lattices model crystals, error-correcting codes, cryptography, optimization, and integer approximation.

\term{Laurent Series}
A Laurent series represents a complex function near $z_0$ as $f(z)=\sum_{n=-\infty}^{\infty}a_n(z-z_0)^n$. Negative powers describe the singular part: finitely many indicate a pole, while infinitely many may indicate an essential singularity. \par\noindent\textbf{Applications:} Laurent series calculate residues and evaluate contour integrals.

\term{L-function}
An $L$-function is a Dirichlet series, often with an Euler product, that generalizes the Riemann zeta function; for a character $\chi$, $L(s,\chi)=\sum\chi(n)n^{-s}$. Its analytic behavior encodes arithmetic information. \par\noindent\textbf{Applications:} $L$-functions study primes, elliptic curves, modular forms, and number-field extensions.

\term{L-space}
In Heegaard Floer homology, an L-space is a 3-manifold $Y$ whose Heegaard Floer homology $\widehat{HF}(Y)$ is as small as possible (rank equal to the order of $H_1(Y; \mathbb{Z})$).

\term{Law of Large Numbers}
The law of large numbers says that averages of many independent identically distributed variables approach their common expectation. The weak law gives convergence in probability, while the strong law gives almost-sure convergence under suitable conditions. \par\noindent\textbf{Applications:} it justifies sample averages, Monte Carlo computation, polling, insurance, and statistical estimation.

\term{Law of Total Probability}
If $\{A_i\}$ is a partition, then $P(B)=\sum_iP(B\mid A_i)P(A_i)$. It decomposes an event according to mutually exclusive cases. \par\noindent\textbf{Applications:} it supports Bayesian inference, reliability calculations, risk analysis, and prediction with hidden categories.

\term{Least Common Multiple}
The least common multiple $\operatorname{lcm}(a,b)$ is the smallest positive integer divisible by both $a$ and $b$. For positive integers, $\gcd(a,b)\operatorname{lcm}(a,b)=ab$. \par\noindent\textbf{Applications:} LCM handles synchronized cycles, common denominators, modular schedules, and periodic processes.

\term{Lax Equivalence Theorem}
For a consistent finite-difference scheme approximating a well-posed linear initial-value problem, stability is equivalent to convergence. Consistency controls local truncation error, while stability prevents errors from growing without bound. \par\noindent\textbf{Applications:} the theorem guides reliable numerical schemes for wave, transport, and other evolution equations.

\term{Lax--Milgram Theorem}
If $a(u,v)$ is bounded and coercive on a Hilbert space $H$, then for every bounded linear functional $f$ there is a unique $u\in H$ with $a(u,v)=f(v)$ for all $v$. The solution depends continuously on $f$. \par\noindent\textbf{Applications:} it proves existence and stability of weak solutions to elliptic PDEs and supports finite-element methods.

\term{Least Squares}
Least squares chooses parameters to minimize the sum of squared residuals between observations and model predictions. For a linear model this leads to the normal equations, or a more stable QR or singular-value computation. \par\noindent\textbf{Applications:} least squares is used in regression, calibration, signal reconstruction, inverse problems, and data fitting.

\term{Lebesgue Integral}
The Lebesgue integral groups points according to function values and measures the corresponding sets, rather than relying only on intervals in the domain. It extends the Riemann integral and handles limits of functions and many discontinuities effectively. \par\noindent\textbf{Applications:} it is the foundation of modern probability, Fourier analysis, PDEs, and $L^p$ spaces.

\term{Lebesgue Measure}
Lebesgue measure assigns length, area, or volume to measurable subsets of $\mathbb R^n$. It is translation invariant and agrees with ordinary volume on boxes, with the unit cube having measure $1$. \par\noindent\textbf{Applications:} it defines probability densities, integration, volume in analysis, and the size of exceptional or fractal sets.

\term{Lebesgue's Dominated Convergence Theorem}
If measurable $f_n\to f$ almost everywhere and $|f_n|\leq g$ for an integrable $g$, then $\int f_n\to\int f$. The common integrable bound permits the limit to pass through the integral. \par\noindent\textbf{Applications:} it justifies limiting calculations in probability, Fourier analysis, numerical approximation, and PDEs.

\term{Lusin's Theorem}
On a finite-measure space, a measurable function is nearly continuous: for every $\varepsilon>0$, there is a closed set whose complement has measure below $\varepsilon$ and on which the function is continuous. \par\noindent\textbf{Applications:} Lusin's theorem connects measurable functions with continuous approximation and supports density results in $L^p$ analysis.

\term{Lemma}
A "helping theorem"; a minor result that is proved as a stepping stone toward a larger, more significant theorem.

\term{Leibniz Rule (Generalised)}
The generalized Leibniz rule is
\[
(fg)^{(n)}=\sum_{k=0}^n\binom nk f^{(k)}g^{(n-k)}.
\]
It extends the product rule to higher derivatives. \par\noindent\textbf{Applications:} it is used in Taylor expansions, differential equations, symbolic computation, and estimates for products of smooth functions.

\term{Leibniz Integral Rule}
The Leibniz rule differentiates an integral with variable limits:
\[
\frac d{dx}\int_{a(x)}^{b(x)}f(x,t)dt=f(x,b(x))b'(x)-f(x,a(x))a'(x)+\int_{a(x)}^{b(x)}f_x(x,t)dt.
\]
\par\noindent\textbf{Applications:} it derives conservation laws, sensitivity equations, moving-boundary models, and gradients of integral objectives.

\term{Levi--Civita Connection}
The Levi--Civita connection is the unique connection that is torsion-free and preserves the Riemannian metric. It defines covariant differentiation, parallel transport, curvature, and geodesics. \par\noindent\textbf{Applications:} it provides the basic language for curved surfaces, general relativity, mechanics, and geometric data analysis.

\term{L'Hôpital's Rule}
Under appropriate differentiability and nonzero-denominator hypotheses, an indeterminate limit of $f/g$ of type $0/0$ or $\infty/\infty$ may be found from $f'/g'$ when that latter limit exists. \par\noindent\textbf{Applications:} it evaluates difficult limits in calculus and asymptotic error analysis; its hypotheses must be checked before use.

\term{Liar Paradox}
The self-referential sentence "this statement is false" cannot consistently be assigned a truth value, since it is true iff it is false. It is a foundational example of self-reference in logic, motivating Tarski's hierarchy of truth and discussions of paraconsistent and gap-based approaches.

\term{Lie Algebra}
A Lie algebra is a vector space with a bilinear bracket satisfying $[X,X]=0$ and the Jacobi identity. The bracket measures the failure of infinitesimal operations to commute. \par\noindent\textbf{Applications:} Lie algebras linearize Lie groups and describe rotations, differential equations, symmetries, gauge fields, and conservation laws.

\term{Lie Bracket}
The Lie bracket measures noncommutativity. For matrices or operators it is $[X,Y]=XY-YX$; for vector fields it records their commutator as differential operators. It is bilinear, antisymmetric, and satisfies Jacobi. \par\noindent\textbf{Applications:} brackets describe infinitesimal symmetries, controllability, curvature, and quantum commutators.

\term{Lie Group}
A Lie group is both a group and a smooth manifold, with smooth multiplication and inversion. Its tangent space at the identity is a Lie algebra describing infinitesimal motions. \par\noindent\textbf{Applications:} Lie groups model continuous symmetries in robotics, mechanics, computer vision, relativity, and physics.

\term{Lie Derivative}
The Lie derivative measures how a tensor field changes when carried along the flow of a vector field. For a scalar it is directional differentiation; for vector fields it is related to the Lie bracket. \par\noindent\textbf{Applications:} it tests symmetries, expresses conservation laws, and is used in fluid mechanics, differential geometry, and relativity.

\term{Lie Superalgebra}
A generalization of a Lie algebra that is $\mathbb{Z}_2$-graded (even and odd parts) and satisfies a graded version of the Jacobi identity. Used in supersymmetry in physics.

\term{Limit}
The limit of a function or sequence is the value it approaches as its input approaches a point or infinity. The $\varepsilon$--$\delta$ definition makes this idea precise without requiring the function to attain the limiting value. \par\noindent\textbf{Applications:} limits define derivatives, integrals, continuity, convergence, and asymptotic approximations.

\term{Limit (Category Theory)}
A limit of a diagram $F: J \to C$ is a cone $(L, \pi_j)$ over $F$ that is universal: every other cone factors uniquely through $L$. Limits encompass products, pullbacks, equalizers, and inverse limits. A category with all small limits is complete, and limits are computed objectwise in functor categories. Right adjoints preserve limits, and their existence is governed by the adjoint functor theorems. Limits are the dual of colimits and central to universal constructions; this is distinct from the analytic notion of a limit in calculus.

\term{Limit Point}
A point $x$ is a limit or accumulation point of a set $S$ if every neighborhood of $x$ contains a point of $S$ different from $x$. In metric spaces this is equivalent to being the limit of a sequence of distinct points of $S$. \par\noindent\textbf{Applications:} limit points describe closure, continuity, convergence, and the boundary structure of sets.

\term{Lindenbaum--Tarski Algebra}
The algebra formed by the equivalence classes of formulas in a logical theory, where the equivalence relation is provable equivalence.

\term{Lindeberg--Lévy Central Limit Theorem}
The Lindeberg--Lévy central limit theorem states that normalized sums of independent identically distributed variables with finite mean and nonzero finite variance converge in distribution to a standard normal variable. \par\noindent\textbf{Applications:} it explains normal approximations for sample means, measurement errors, polling, and statistical confidence intervals.

\term{Lindelöf Space}
A Lindelöf space is one in which every open cover has a countable subcover. This property provides countable control over potentially large collections of open sets; every second-countable space is Lindelöf. \par\noindent\textbf{Applications:} it supports separability and countable constructions in analysis and topology.

\term{Line}
A one-dimensional object with zero thickness extending without end in both directions; in a projective plane, the dual of a point.

\term{Line Bundle}
A line bundle is a vector bundle with one-dimensional fibers. Local sections look like functions whose descriptions change by nonzero scaling on overlaps. Divisors and line bundles are closely related, and the first Chern class records topological twisting. \par\noindent\textbf{Applications:} line bundles define projective embeddings, describe phases and magnetic potentials, and organize divisors in algebraic and complex geometry.

\term{Line Integral}
A line integral accumulates a scalar field along arc length, $\int_Cf\,ds$, or measures work and circulation of a vector field, $\int_C\mathbf F\cdot d\mathbf r$. Its value can depend on the path unless the field is conservative. \par\noindent\textbf{Applications:} line integrals compute work, circulation, mass of wires, and quantities along trajectories.

\term{Linear Algebra}
Linear algebra studies vector spaces, linear maps, matrices, and linear equations. Concepts such as bases, rank, eigenvalues, and orthogonality turn complicated systems into structured calculations. \par\noindent\textbf{Applications:} it is fundamental in engineering, data science, graphics, quantum mechanics, differential equations, and optimization.

\term{Linear Combination}
A linear combination of vectors is $\sum_i a_iv_i$, where the $a_i$ are scalars. The span of a set is all its linear combinations. \par\noindent\textbf{Applications:} linear combinations describe coordinates, approximations, mixtures, signals, and solutions of linear systems.

\term{Linear Dependence}
A set of vectors $\{v_i\}$ is linearly dependent if there exist scalars $a_i$, not all zero, such that $\sum a_i v_i = 0$. Otherwise, they are linearly independent.

\term{Linear Functional}
A linear functional is a linear map $\ell:V\to\mathbb F$, satisfying $\ell(au+bv)=a\ell(u)+b\ell(v)$. All such maps form the dual space $V^*$. \par\noindent\textbf{Applications:} functionals represent measurements, constraints, gradients, weak derivatives, and observations in optimization and PDEs.

\term{Linear Independence}
A set of vectors is linearly independent if the only linear combination equal to zero has all coefficients zero.

\term{Linear Map}
A linear map satisfies $T(au+bv)=aT(u)+bT(v)$. Once bases are chosen, it is represented by a matrix, so kernels, images, rank, and eigenvalues can be studied algebraically. \par\noindent\textbf{Applications:} linear maps model transformations, networks, discretized physical systems, and changes of coordinates.

\term{Linear Programming}
Linear programming optimizes a linear objective subject to linear equalities and inequalities. The feasible set is a convex polyhedron, and when an optimum exists it can be found at an extreme point. \par\noindent\textbf{Applications:} it allocates resources, schedules production, routes goods, designs networks, and supports machine learning.

\term{Linear Recurrence}
A linear recurrence expresses each term as a fixed linear combination of earlier terms, such as $a_n=c_1a_{n-1}+\cdots+c_ka_{n-k}$. Its characteristic polynomial often gives a closed form. \par\noindent\textbf{Applications:} recurrences model population growth, algorithms, digital filters, and discrete dynamical systems.

\term{Linear Regression}
Linear regression models a response as $y\approx\beta_0+\beta_1x_1+\cdots+\beta_px_p+\varepsilon$. Parameters are often estimated by least squares, with assumptions used for uncertainty and inference. \par\noindent\textbf{Applications:} regression supports prediction, calibration, trend estimation, and experimental analysis.

\term{Linear Operator}
A linear map from a vector space to itself. The study of these operators includes eigenvalues, eigenvectors, and spectral theory.

\term{Lipschitz Continuity}
A function is Lipschitz continuous if $|f(x)-f(y)|\leq L|x-y|$ for all $x,y$. It cannot change faster than a fixed linear rate. Lipschitz continuity implies uniform continuity. \par\noindent\textbf{Applications:} it guarantees uniqueness of ODE solutions under standard hypotheses and controls approximation error, stability, and learning models.

\term{Lipschitz Constant}
The Lipschitz constant is the smallest admissible $L$ in $|f(x)-f(y)|\leq L|x-y|$. It bounds the maximum global rate of change. \par\noindent\textbf{Applications:} it quantifies sensitivity, controls propagation of errors, and sets step-size or robustness bounds in numerical methods and optimization.

\term{Liouville's Theorem}
Liouville's theorem states that every bounded entire holomorphic function is constant. Cauchy's estimates force all positive-degree Taylor coefficients to vanish. \par\noindent\textbf{Applications:} it proves the fundamental theorem of algebra and gives powerful uniqueness and growth arguments in complex analysis.

\term{Local Compactness}
A space is locally compact if every point has a neighborhood whose closure is compact, or equivalently a compact neighborhood in common Hausdorff settings. Local compactness allows useful integration and compactness arguments without the whole space being compact. \par\noindent\textbf{Applications:} it underlies Haar measure, Fourier analysis, potential theory, and one-point compactification.

\term{Local Field}
A locally compact, non-discrete topological field: the real or complex numbers, a finite extension of the field $\mathbb{Q}_p$ of $p$-adic numbers, or a finite extension of the field $\mathbb{F}_q((t))$ of formal Laurent series over a finite field. Every non-archimedean local field carries a discrete valuation and a finite residue field. Local class field theory describes abelian extensions of local fields, Tate duality governs local Galois cohomology, and Haar measure on the additive group provides the local structure underlying the adèles.

\term{Local Homeomorphism}
A map is a local homeomorphism if every point has a neighborhood on which the map is a homeomorphism onto an open subset of the target. It need not be globally one-to-one. \par\noindent\textbf{Applications:} local homeomorphisms describe covering maps, coordinate charts, and locally invertible transformations.

\term{Local Minimum}
A point $x_0$ is a local minimum if $f(x_0)\leq f(x)$ for all $x$ in some neighborhood of $x_0$. It is a local comparison and may not be the smallest value globally. \par\noindent\textbf{Applications:} local minima represent stable designs and equilibria and are targets of many optimization algorithms.

\term{Local Maximum}
A point $x_0$ is a local maximum if $f(x)\leq f(x_0)$ for all $x$ in some neighborhood of $x_0$. It need not be the largest value on the whole domain. \par\noindent\textbf{Applications:} local maxima describe nearby best responses, peaks in data, and stable or unstable equilibria.

\term{Local Ring}
A local ring has exactly one maximal ideal, representing the functions that vanish at the point under study. In algebraic geometry, localizing at a prime ideal focuses on behavior near that point. \par\noindent\textbf{Applications:} local rings study singularities, tangent spaces, divisors, and local solvability of equations.

\term{Localization}
The process of formally adjoining inverses of a multiplicatively closed subset $S$ of a commutative ring $R$, producing the ring of fractions $S^{-1}R$; localization at the complement of a prime ideal $\mathfrak{p}$ yields the local ring $R_{\mathfrak{p}}$. For modules, localization is exact and compatible with tensor products. Geometrically, localization corresponds to restricting sheaves to an open set, expressing the local behaviour of a variety at a point. The notion generalizes to the localization of a category, formally inverting a class of morphisms, as in the construction of derived categories.

\term{Logarithm}
The logarithm is the inverse of exponentiation: $\log_bx=y$ when $b^y=x$. The natural logarithm satisfies $\ln x=\int_1^xdt/t$ and converts products into sums. \par\noindent\textbf{Applications:} logarithms measure growth, scale data, solve exponential equations, and appear in entropy and decibel units.

\term{Logarithmic Function}
The inverse of the exponential function, $\log_b x = y$ when $b^y = x$, satisfying $\log(xy) = \log x + \log y$.

\term{Loop Space}
The loop space $\Omega X$ consists of paths in a pointed space that start and end at the basepoint. Concatenation gives it a group-like structure up to homotopy, and $\pi_n(X)\cong\pi_{n-1}(\Omega X)$. \par\noindent\textbf{Applications:} loop spaces convert questions about maps and homotopy groups into structured algebraic-topology problems and are central to stable homotopy theory.

\term{Lorentz Transformation}
A Lorentz transformation is a linear change between inertial frames preserving the spacetime interval $\Delta s^2=-c^2\Delta t^2+\Delta x^2+\Delta y^2+\Delta z^2$. It combines rotations with boosts and replaces ordinary Euclidean distance by spacetime geometry. \par\noindent\textbf{Applications:} Lorentz transformations describe relativistic motion, time dilation, length contraction, and electromagnetic-field changes.

\term{Löwenheim--Skolem Theorem}
A countable first-order theory that has an infinite model has a model of any infinite cardinality (downward and upward forms). It yields the Skolem paradox: set theory has a countable model, yet proves the existence of uncountable sets.

\term{Lower Bound}
A lower bound of a set $S$ is a number $m$ with $m\leq s$ for every $s\in S$. The greatest lower bound is $\inf S$, whether or not it belongs to $S$. \par\noindent\textbf{Applications:} lower bounds certify performance limits, prove convergence, and establish minimum possible costs or errors.

\term{L\textasciicircum p Space}
The space $L^p$ contains measurable functions with $\int|f|^p<\infty$ for $1\leq p<\infty$, with an essential-supremum norm when $p=\infty$. For $p=2$ it is a Hilbert space. \par\noindent\textbf{Applications:} $L^p$ spaces measure signal size, define probability moments, and provide the standard setting for Fourier analysis and PDEs.

\term{L-polynomial}
The polynomial associated with an L-function, often appearing in the context of the Birch and Swinnerton-Dyer conjecture for elliptic curves.

\term{Lefschetz Fixed-Point Theorem}
The Lefschetz fixed-point theorem relates fixed points of a continuous map to its action on homology. The Lefschetz number is $L(f)=\sum_k(-1)^k\operatorname{tr}(f_*|H_k(X))$; if it is nonzero, $f$ has a fixed point. \par\noindent\textbf{Applications:} it proves existence of equilibria and fixed points in topology, dynamics, and geometry.

\term{Lefschetz Hyperplane Theorem}
A result in complex geometry stating that the cohomology of a smooth projective variety $X$ is determined by the cohomology of a generic hyperplane section $H$ up to a certain dimension.

\term{L-system (Lindenmayer system)}
An L-system is a parallel rewriting system: every symbol in a string is replaced simultaneously according to production rules. Interpreting symbols as drawing commands produces branching patterns. \par\noindent\textbf{Applications:} L-systems model plant growth, generate fractals, and create procedural graphics.

\term{Lattice Gauge Theory}
Lattice gauge theory approximates spacetime by a discrete lattice while preserving gauge symmetry through variables on edges. It provides a non-perturbative formulation of quantum field theories. \par\noindent\textbf{Applications:} numerical lattice calculations study strong interactions, confinement, phase transitions, and quantum statistical systems.

\term{Lumpability}
Lumpability is a property of a Markov chain under which grouped states form an exact lower-dimensional Markov chain. Transition probabilities from states in the same group must agree after aggregation. \par\noindent\textbf{Applications:} lumpability enables model reduction, state aggregation, faster simulation, and simplified queueing or reliability models.
